% merged_thesis.tex — formal thesis wrapper using merged chapters
% Compile from project root:
%   pdflatex thesis_revised_merged/merged_thesis && bibtex merged_thesis && pdflatex thesis_revised_merged/merged_thesis x2
% Or copy this file to project root and adjust \input paths if your build expects root-level .tex.

\documentclass[12pt]{report}

\usepackage[language=english, type=bachelor]{style}
\usepackage{amsmath}
\usepackage{amssymb}

\newcommand{\thesisfig}[2]{%
  \IfFileExists{#2}{%
    \includegraphics[width=#1]{#2}%
  }{%
    \fbox{\parbox{#1}{\centering\vspace{1.1cm}\footnotesize\textit{Add file:} \texttt{\detokenize{#2}}\vspace{1.1cm}}}%
  }%
}

\usepackage{fancyhdr}
\setlength{\headheight}{27.35799pt}
\pagestyle{fancy}
\lhead{}
\chead{}
\renewcommand{\headrulewidth}{0.2pt}
\renewcommand{\footrulewidth}{0.2pt}

\begin{document}

\specialization{Computer Science}
\title{Augmented Reality-Based Interior Design: A Mobile Application for Parametric Room Scanning and Interactive 3D Scene Reconstruction}
\author{Mara Ielciu}
\supervisor{Professor Dr. Darabant Sergiu Adrian}

\maketitle

\newpage
\thispagestyle{empty}
\mbox{}
\newpage
\pagenumbering{roman}

\cleardoublepage
ABSTRACT
\vspace{0.5cm}
\hrule
\vspace{0.5cm}

This bachelor thesis presents an iOS application for augmented reality--based interior design. Users scan a real room with Apple's RoomPlan API to obtain a parametric, semantically labelled model (\texttt{CapturedRoom}: walls, doors, windows, and detected furniture with metric dimensions). The scan is loaded into a non-AR three-dimensional editor built with RealityKit, where detected furniture can be replaced with USDZ models from a library or produced by Object Capture's on-device photogrammetry pipeline. Edits persist locally (JSON scan data under the app Documents directory, Codable room models). The application computes room measurements and renovation estimates from scan geometry, provides rule-based interior-design suggestions, and exports completed designs as glTF-based ZIP archives with USDZ furniture and a Blender helper script for further editing in desktop DCC tools. For visual comparison, the project supports importing pretrained 3D Gaussian Splatting assets (\texttt{.ply}) and requesting feedforward SHARP reconstructions from a user-controlled Mac on the local network, rendering both through a custom Metal splat viewer. A three-way comparison among mesh-based staging, offline-trained splats, and SHARP clarifies trade-offs among geometric editability, capture effort, and plenoptic fidelity.

Chapter~2 reviews photogrammetry (structure from motion, multi-view stereo, Poisson surface reconstruction), depth fusion via truncated signed distance functions, and how Object Capture versus RoomPlan map onto image-based versus depth-assisted reconstruction. Chapter~3 summarises 3D Gaussian Splatting, feedforward SHARP inference, and the mesh-versus-splat design choice. Chapters~4 and~5 state requirements and describe the software architecture (SwiftUI, MVVM-oriented structure, major views and services), data flow, analytics, Blender export, and testing. The main contribution is an integrated, on-device workflow from room scan to persistent, editable virtual staging, with optional radiance-field visualisation and professional hand-off via Blender export.

\tableofcontents

\newpage
\pagenumbering{arabic}

\section*{Context and Related Work}
\addcontentsline{toc}{section}{Context and Related Work}

Indoor scene understanding and 3D reconstruction have been central to computer vision research, with datasets and benchmarks driving progress in semantic segmentation and object detection. ScanNet~\cite{dai2017scannet} is a large-scale dataset of richly annotated 3D reconstructions of indoor scenes, with 2.5 million views across more than 1,500 scans and instance-level semantic segmentations. That work established the importance of combining geometric reconstruction with semantic labelling (walls, furniture, openings) for applications such as robotics and spatial reasoning. The present application aligns with this direction by leveraging Apple's RoomPlan to produce a parametric, semantically annotated room model from a consumer device, rather than from offline processing of RGB-D videos.

The use of augmented reality for interior design and furniture visualisation has been explored in both industry and academia. K\'an~et~al.~\cite{kan2021ar} present an AR system for automatic interior design based on a hierarchical tree of procedural rules, enabling interactive and personalised furniture configuration with a recommender component; they compare AR-based visualisation to traditional methods and report positive user feedback. Such studies confirm that AR can support decision-making and spatial reasoning when users need to envision furniture and finishes in their own spaces. This thesis extends that idea by allowing users not only to place virtual furniture, but to \emph{replace} detected real furniture in a reconstructed 3D scene, thus supporting ``what-if'' redesign directly tied to the current room layout.

User experience in mobile AR interior design applications has been studied to improve usability and adoption. Kandil~et~al.~\cite{kandil2021ux} evaluate interior design mobile AR applications using established scales such as the System Usability Scale (SUS) and the Handheld Augmented Reality Usability Scale (HARUS), with a focus on gesture-based interaction and one-handed use. Their findings underline the need for clear onboarding, stable tracking, and intuitive manipulation of virtual objects. The application developed in this thesis addresses these concerns by using RoomPlan's guided scanning flow, RealityKit's responsive rendering, and SwiftUI's declarative interface to keep interactions consistent and learnable.

\chapter{Introduction}
\label{ch:introduction}

\section{Motivation and Objectives}
\label{sec:intro:motivation}

The problem addressed by this thesis arises at the intersection of two
well-documented shortcomings in current practice. On one side, conventional
interior design tools, whether two-dimensional floor plans, printed catalogues,
or digital product configurators, present furniture and spatial layouts as
abstractions that strip away the perceptual properties that ultimately determine
whether a redesign succeeds: how an object occludes a window, whether a sofa
visually dominates a narrow corridor, or how a chosen material interacts with
the room's natural light. On the other side, consumer augmented reality
applications that \emph{do} allow virtual objects to be previewed in a live
camera feed share a structural limitation identified empirically by
Kandil~et~al.~\cite{kandil2021ux}: they treat the room as a transiently
detected floor plane, provide no persistent room model, and require the user
to be physically present in the space for every design iteration. The result
is a tool that is useful for a single in-store moment but not for the extended,
iterative process that characterises real interior redesign decisions.

A more principled approach is to treat the room scan itself as the primary
artefact of the workflow. If the real room is captured once as a metric,
semantically annotated three-dimensional model, a \emph{digital twin}, subsequent
design iterations can be performed at any time and from any location without
requiring the user to return to the physical space. Indoor scene understanding
research has demonstrated that this level of structure is achievable at scale
on consumer hardware: ScanNet~\cite{dai2017scannet} showed that richly annotated
RGB-D reconstructions with per-instance semantic labels support a wide range of
downstream tasks, and Matterport3D~\cite{CAS+17} extended that work to
building-scale scenes with 40-category instance labels. Apple's RoomPlan
API~\cite{apple2024roomplan2} brings a real-time, on-device subset of this
capability to supported iPhones and iPads, outputting a parametric
\texttt{CapturedRoom} whose typed structural surfaces and classified furniture
instances reflect the same conceptual taxonomy that these research datasets
made trainable.

A second motivation concerns the asset authoring bottleneck. Replacing a
detected furniture item with a virtual equivalent requires a three-dimensional
model of that item. Static libraries cover a limited catalogue and exclude
bespoke or non-commercial pieces. Object Capture~\cite{apple2024roomplan2}
enables the user to scan a physical object with the same device and obtain a
USDZ mesh through an on-device photogrammetric pipeline, closing the loop
between room capture and asset creation without external tooling during the
core staging workflow.

A third motivation concerns representation choice and workflow completeness.
3D Gaussian Splatting~\cite{kerbl20233dgs} achieves real-time rendering quality
competitive with neural radiance fields while remaining differentiable for
offline training. In this application, explicit meshes remain the medium for
interactive furniture replacement in RealityKit, whereas imported splats and
feedforward SHARP reconstructions serve as \emph{visual benchmarks}. Users who
require desktop refinement can export staged rooms to Blender as glTF-based ZIP
archives; users who need rapid radiance-field previews can invoke a local SHARP
inference server. Chapter~\ref{ch:theory2} formalises the three object-level
reconstruction paths (Assets~A, B, and~C); Chapter~\ref{ch:implementation}
documents their implementation.

The work pursues the following objectives, each directly justified by the
gaps identified above:

\begin{enumerate}
  \item \textbf{Room-scale semantic capture.}
    Integrate RoomPlan so that a scan yields a \texttt{CapturedRoom} containing
    typed structural elements (walls, floors, ceilings, doors, windows) and
    classified furniture instances with metric transforms and bounding dimensions.
    \emph{Justification:} without a persistent, semantically structured room
    model, furniture replacement cannot be grounded to real geometry.

  \item \textbf{Persistent virtual staging.}
    Provide a non-AR three-dimensional editor with orbit-style navigation and
    gesture-based manipulation (translation, rotation, scale), with state
    persisted across application launches.
    \emph{Justification:} Kandil~et~al.~\cite{kandil2021ux} found that
    requiring physical presence and absence of session persistence were primary
    usability barriers in existing AR interior apps.

  \item \textbf{Object-scale asset authoring (Asset~A).}
    Integrate Object Capture for guided multi-view photogrammetric
    reconstruction into USDZ, extending the furniture library beyond static
    catalogue content.

  \item \textbf{Radiance-field benchmarks (Assets~B and~C).}
    Import offline-trained 3DGS \texttt{.ply} files (Asset~B) and invoke Apple's
    SHARP feedforward model on a LAN server from a single photograph (Asset~C);
    render both through a shared Metal splat viewer for qualitative comparison
    with meshes---not for furniture editing.

  \item \textbf{Room analytics and planning aids.}
    Compute metric dimensions and renovation estimates from scan geometry;
    provide rule-based interior-design suggestions (including lighting density
    heuristics) in a dedicated statistics interface.

  \item \textbf{DCC hand-off.}
    Export completed rooms as glTF-based ZIP archives with USDZ furniture,
    manifests, and a Blender helper script for further editing outside the device.
\end{enumerate}


\section{Contributions and Scope}
\label{sec:intro:contributions}

The contributions of this thesis are primarily integrative and
engineering-oriented, consistent with the scope of a bachelor project:

\begin{itemize}
  \item \textbf{A unified scan-to-edit pipeline.}
    A complete iOS application connecting RoomPlan capture to a serialised
    room data model, to a RealityKit scene graph, and to furniture replacement
    with persistence across application launches.

  \item \textbf{Theoretical framing of Apple's capture APIs within the
    academic literature.}
    A structured account of how Object Capture instantiates the classical
    photogrammetric chain and how RoomPlan builds on depth fusion in the TSDF
    family~\cite{newcombe2011kinectfusion}, augmented by learned structure and
    semantics, with a principled comparison to the 3D Gaussian Splatting
    rendering path~\cite{kerbl20233dgs}.

  \item \textbf{A three-way comparative analysis} of mesh (Asset~A),
    offline 3DGS (Asset~B), and SHARP splats (Asset~C), clarifying roles:
    interactive staging versus visual benchmarking.

  \item \textbf{A Metal-based Gaussian splat renderer for iOS.}
    Loading standard binary PLY files from common training frameworks at
    interactive frame rates, bridged into SwiftUI via \texttt{MTKView}.

  \item \textbf{Supporting services.}
    Two-level \texttt{CapturedRoom} caching, rule-based
    \texttt{SmartSuggestionsService}, scan-time colour sampling, Blender export,
    and local-network SHARP integration without third-party cloud upload for
    Apple capture paths.
\end{itemize}

The thesis focuses on software architecture, core algorithms at a conceptual
level, and the iOS platform. It does not contribute novel learning
architectures for room understanding: RoomPlan's neural models are Apple's
proprietary intellectual property. 3DGS training is assumed to run off-device
in the current system; on-device splat training remains
research-grade~\cite{Guo26}. Full LiDAR-quality RoomPlan output requires
iPhone~12~Pro or later, or a LiDAR-equipped iPad~Pro; the application
degrades gracefully on non-LiDAR devices where Apple's API permits.

\section{Document Structure}

Chapter~\ref{ch:theory1} presents photogrammetry, depth fusion, and the mapping to Object Capture and RoomPlan, followed by related work. Chapter~\ref{ch:theory2} introduces 3D Gaussian Splatting, feedforward SHARP inference, the mesh-versus-splat design choice, and the three-asset comparison. Chapter~\ref{ch:requirements} formalises functional and non-functional requirements and system constraints. Chapter~\ref{ch:implementation} describes architecture, data flow, implementation details, analytics, Blender export, the user interface, and testing. Chapter~\ref{ch:conclusions} summarises results, limitations, and future work.

The opening \textit{Context and Related Work} section situates the project against ScanNet-style indoor semantics~\cite{dai2017scannet}, AR interior design systems~\cite{kan2021ar}, and mobile AR usability evaluation~\cite{kandil2021ux}.

\chapter{Photogrammetry, Depth Fusion, and Apple Capture APIs}
\label{ch:theory1}

This chapter explains the geometric and algorithmic ideas behind the capture pipelines used in the thesis application. It first reviews \textbf{photogrammetry} and the classical chain \textbf{structure from motion (SfM) $\rightarrow$ multi-view stereo (MVS) $\rightarrow$ Poisson surface reconstruction} \cite{luhmann2014photogrammetry,lowe2004sift,furukawa2010pmvs,kazhdan2006poisson}. It then describes \textbf{depth fusion} with a \textbf{truncated signed distance function (TSDF)} as used in real-time RGB-D and LiDAR systems \cite{newcombe2011kinectfusion,whelan2016elasticfusion}, which underpins ARKit-style scene meshes that RoomPlan builds upon. A comparison between \textbf{image-only} object reconstruction and \textbf{depth-assisted} room understanding clarifies why Object Capture and RoomPlan behave differently. Finally, the chapter maps these ideas to Apple's \textbf{Object Capture} and \textbf{RoomPlan} APIs, documents the \texttt{CapturedRoom} data model, and summarises limitations including external refinement via Blender export.

\section{Motivation}
\label{sec:motivation}

Users redesign interiors more confidently when they can \textbf{see} alternatives in context. Traditional representations (2D plans, catalog photos, or material swatches) omit depth, occlusion, and the relation between existing objects and new ones. Augmented reality can composite virtual furniture into a live video stream, but that mode ties the experience to being physically present and holding the device. If the room is first captured as a structured 3D model, the user can reopen and edit that model later: the scan becomes a \textbf{persistent spatial substrate} for experimentation.

Research datasets such as ScanNet~\cite{dai2017scannet} demonstrated that indoor applications benefit from joint \textbf{geometry} and \textbf{semantics} (which surfaces are walls, which objects are chairs or tables). RoomPlan brings a consumer subset of that idea to supported iPhones and iPads: it outputs not only geometry but labelled entities with rough dimensions. This thesis uses that output as the basis for \textbf{furniture replacement}, substituting USDZ models for detected instances, which goes beyond merely placing a free-floating object on a manually identified floor plane.

\section{Photogrammetry and the Classical Reconstruction Pipeline}
\label{sec:photogrammetry}

\subsection{Definition and triangulation}

\textbf{Photogrammetry} is the discipline of extracting metric information from photographs \cite{luhmann2014photogrammetry}. In 3D reconstruction, the elementary geometric operation is \textbf{triangulation}: corresponding points in two or more calibrated images determine a 3D point by intersecting rays. Stable triangulation requires \textbf{baseline} between cameras and sufficient \textbf{overlap} so that the same surface points appear in multiple views.

\begin{figure}[htbp]
\centering
\thesisfig{0.82\textwidth}{figures/ch2_photogrammetry_concept.png}
\caption{Conceptual photogrammetry model: multi-view image observations, camera parameters, and additional constraints used to estimate 3D coordinates.}
\label{fig:ch2_photo_concept}
\end{figure}

\subsection{From features to textured mesh}

A widely used multi-view pipeline proceeds in stages \cite{luhmann2014photogrammetry,lowe2004sift,furukawa2010pmvs,kazhdan2006poisson}:

\begin{enumerate}
    \item \textbf{Feature detection and matching}: Distinctive keypoints (e.g.\ SIFT \cite{lowe2004sift}) are extracted and matched across images.

    \item \textbf{Structure from Motion}: Camera intrinsics and extrinsics are estimated together with a \emph{sparse} 3D point cloud. \textbf{Bundle adjustment} refines parameters by minimising reprojection error across all observations.

    \item \textbf{Multi-view stereo}: Dense correspondence and depth estimation across views yield a dense point cloud or per-view depth maps \cite{furukawa2010pmvs}.

    \item \textbf{Surface reconstruction}: The dense samples are converted to a watertight or oriented surface, commonly using \textbf{Poisson surface reconstruction} \cite{kazhdan2006poisson}.

    \item \textbf{Texture mapping}: The mesh is parameterised; colours are blended from source images onto the surface.
\end{enumerate}

\begin{figure}[htbp]
\centering
\thesisfig{0.92\textwidth}{figures/fig2_2_sfm_mvs}
\caption{SfM--MVS reconstruction pipeline: from multi-view image input through feature detection, Structure from Motion, dense Multi-View Stereo, to surface reconstruction and textured mesh output (adapted from \cite{Kholil_2021}).}
\label{fig:ch2_photo_pipeline}
\end{figure}

This end-to-end path is what Section~\ref{sec:on_device_pipeline} relates to \textbf{Object Capture} (Asset~A in Chapter~\ref{ch:theory2}).

\section{Depth Fusion and TSDF-Based Scene Meshing}
\label{sec:depth_fusion}

Room-scale tracking on mobile devices fuses \textbf{depth} (LiDAR or stereo) with \textbf{visual--inertial odometry}. A standard volumetric representation is the \textbf{TSDF}: for each voxel in a bounded volume, one stores a (truncated) signed distance to the nearest surface along with a weight; each new depth frame updates voxels along viewing rays \cite{newcombe2011kinectfusion}. Over time, noise averages down and a coherent surface emerges. ElasticFusion~\cite{whelan2016elasticfusion} extended the idea to room-scale consistency when loop closures occur.

Let $D_t(\mathbf{x})$ denote the signed distance implied by depth frame $t$ at voxel $\mathbf{x}$, with confidence weight $w_t$. A common fusion rule updates the field $F$ and weight $W$ as
\[
  F_{t+1}(\mathbf{x}) = \frac{W_t(\mathbf{x}) F_t(\mathbf{x}) + w_t D_t(\mathbf{x})}{W_t(\mathbf{x}) + w_t},
  \qquad
  W_{t+1}(\mathbf{x}) = W_t(\mathbf{x}) + w_t,
\]
with truncation so that voxels far from the surface stop updating. ARKit's scene mesh used during RoomPlan scanning is, conceptually, an isosurface of such a field, though Apple does not expose the raw voxel grid.

Semantic interpretation (which planar patch is a ``wall'' and which blob is a ``sofa'') is not solved by the TSDF alone; RoomPlan adds learning-based detection and classification on top of ARKit scene understanding.

\section{Image-Only versus Depth-Assisted Reconstruction}
\label{sec:depth_vs_photo}

\textbf{Object Capture} is primarily an \textbf{image-driven} pipeline (with poses from ARKit): geometry is inferred from many RGB views through stages analogous to SfM, MVS, and Poisson-class meshing, followed by texturing. \textbf{RoomPlan} is \textbf{depth- and semantics-first}: it targets a structured, parametric room description (\texttt{CapturedRoom}) rather than a single closed mesh of the whole interior. Table~\ref{tab:depth_vs_photo} contrasts the two paradigms at the level relevant to this thesis.

\begin{table}[htbp]
\centering
\caption{Comparison of image-centric object reconstruction (Object Capture) and depth-assisted parametric room capture (RoomPlan).}
\label{tab:depth_vs_photo}
\begin{tabular}{|p{0.28\textwidth}|p{0.30\textwidth}|p{0.30\textwidth}|}
\hline
\textbf{Aspect} & \textbf{Object Capture} & \textbf{RoomPlan} \\
\hline
Primary sensing & Dense RGB coverage, parallax & LiDAR depth + RGB + IMU \\
\hline
Dominant geometry cue & Photo-consistency (MVS) & TSDF fusion + plane fitting \\
\hline
Typical output & Textured triangle mesh (USDZ) & Parametric walls, openings, objects \\
\hline
Semantics & Optional masking & Built-in categories (e.g.\ chair, table) \\
\hline
Role in application & Furniture library (Asset~A) & Editor substrate \\
\hline
\end{tabular}
\end{table}

\section{On-Device Capture Pipeline}
\label{sec:on_device_pipeline}

Figure~\ref{fig:capture_branches} compares two \textbf{parallel pipelines}. The \textbf{top lane} is the photogrammetric/Object Capture path: multi-view RGB images and poses $\rightarrow$ feature matching $\rightarrow$ SfM $\rightarrow$ MVS $\rightarrow$ surface reconstruction and texturing $\rightarrow$ USDZ output. The \textbf{bottom lane} is the RoomPlan path: LiDAR/RGB/IMU input $\rightarrow$ VIO tracking $\rightarrow$ TSDF-class fusion $\rightarrow$ semantic object/opening detection $\rightarrow$ parametric fitting $\rightarrow$ \texttt{CapturedRoom} output. Apple's internal implementations are proprietary; Table~\ref{tab:pipeline_to_api} is a conceptual mapping of stages to public APIs.

\begin{table}[htbp]
\centering
\caption{Conceptual mapping from classical stages to Apple's on-device APIs.}
\label{tab:pipeline_to_api}
\begin{tabular}{|l|p{0.35\textwidth}|p{0.35\textwidth}|}
\hline
\textbf{Stage} & \textbf{Object Capture} & \textbf{RoomPlan} \\
\hline
Camera trajectory & ARKit VIO; session stores images & ARKit world tracking during scan \\
\hline
Sparse / dense geometry & Photogrammetry engine (SfM/MVS class) & TSDF-class fusion + parametric fitting \\
\hline
Surface model & Textured mesh & Structural elements + object boxes \\
\hline
Serialised output & USDZ file & \texttt{CapturedRoom} (Codable) \\
\hline
\end{tabular}
\end{table}

\begin{figure}[H]
\centering
\thesisfig{1\textwidth}{figures/fig2_3_dual_pipeline.pdf}
\caption{Parallel capture pipelines: Object Capture follows an image-based reconstruction chain (top lane); RoomPlan combines LiDAR depth fusion with parametric and semantic inference to produce a \texttt{CapturedRoom} (bottom lane).}
\label{fig:capture_branches}
\end{figure}

\subsection{Object Capture}

\texttt{ObjectCaptureSession} guides the user around an object so that views overlap; \texttt{PhotogrammetrySession} reconstructs a mesh and materials, outputting \textbf{USDZ}. That format is well supported in RealityKit and matches the furniture assets used in the editor and in Blender export copies.

\subsection{Room Plan and \texttt{CapturedRoom}}

Room Plan exposes types such as walls, doors, windows, and \texttt{CapturedRoom.Object} instances with \texttt{category}, \texttt{dimensions}, and \texttt{transform}. Because \texttt{CapturedRoom} conforms to \texttt{Codable}, applications can persist scans as JSON, as done in the implementation (Chapter~\ref{ch:implementation}) under a per-room file in the Documents hierarchy. Scan-time colour sampling attaches optional \texttt{RoomMaterials} to the \texttt{Room} aggregate for editor defaults.

\subsection{Feature-to-API mapping}

Table~\ref{tab:feature_api} ties user-visible capabilities to frameworks.

\begin{table}[h]
\centering
\caption{Major application features and the APIs or artefacts that implement them.}
\label{tab:feature_api}
\begin{tabular}{|l|l|p{0.42\textwidth}|}
\hline
\textbf{Feature} & \textbf{API / artefact} & \textbf{Role} \\
\hline
Room scanning & RoomPlan & Produce \texttt{CapturedRoom} \\
\hline
Room persistence & \texttt{Codable} + file I/O & Reload scans after app restart \\
\hline
3D editor & RealityKit (non-AR) & Visualise and edit room \\
\hline
Object scanning & Object Capture & Author USDZ furniture \\
\hline
Splat viewing & Metal + \texttt{.ply} & Compare with meshes \\
\hline
Blender export & ZIP (glTF + USDZ + script) & DCC hand-off \\
\hline
\end{tabular}
\end{table}

\section{Limitations}
\label{sec:limitations_capture}

Photogrammetry remains sensitive to glass, mirrors, and uniform surfaces. RoomPlan can miss small objects, merge adjacent pieces, or mislabel categories; coarse bounding boxes are not CAD precision. These issues motivate manual touch-up in the editor, \textbf{export to Blender} for desktop refinement of the staged semantic scene (glTF structural geometry, USDZ furniture, and placement manifests), and radiance-field benchmarks (Assets~B and~C) when photographic fidelity matters more than explicit connectivity, as discussed in Chapter~\ref{ch:theory2}. RoomPlan parametric boxes are not a substitute for CAD; Blender export packages the same staging state the editor uses rather than a raw dump of Apple's capture JSON alone.

\chapter{3D Gaussian Splatting}
\label{ch:theory2}

Neural radiance fields (NeRF) showed that a scene can be encoded as a continuous volumetric function approximated by a neural network, enabling high-quality view synthesis from posed images. Training and rendering were initially slow. \textbf{3D Gaussian Splatting} (3DGS), introduced by Kerbl et al.\ \cite{kerbl20233dgs}, replaces the neural field by a set of \textbf{explicit anisotropic Gaussians} in space. Kerbl et al.\ show that this representation delivers \textbf{high visual quality together with real-time rendering}: their method matches or exceeds prior radiance-field quality while running at interactive frame rates on desktop GPUs \cite{kerbl20233dgs}. In this thesis, 3DGS is used as an \textbf{imported} representation (\texttt{.ply} files) rendered with a custom Metal pipeline for \textbf{comparison and viewing}; mesh-based USDZ assets remain the medium for interactive furniture editing and for Blender export. Figure~\ref{fig:ch3_nerf_pipeline} illustrates the NeRF rendering formulation that motivated later explicit real-time representations such as 3DGS.

\begin{figure}[H]
\centering
\thesisfig{0.9\textwidth}{figures/fig3_1_nerf_overview.pdf}
\caption{NeRF overview: sample points along camera rays are fed to a neural MLP predicting colour and density; volume rendering accumulates pixel colours; a photometric loss drives backpropagation (adapted from \cite{mildenhall2020nerf}).}
\label{fig:ch3_nerf_pipeline}
\end{figure}

\section{Representation: 3D Gaussian Primitives}
\label{sec:3dgs_representation}

Each primitive stores a mean $\boldsymbol{\mu} \in \mathbb{R}^3$, a covariance $\boldsymbol{\Sigma}$ describing elongation and orientation, an opacity $\alpha$, and colour parameters. Following Kerbl et al.\ \cite{kerbl20233dgs}, $\boldsymbol{\Sigma}$ is built from a rotation matrix $\mathbf{R}$ (quaternion $\mathbf{q}$) and diagonal scale $\mathbf{S}=\operatorname{diag}(\mathbf{s})$ with positive entries: $\boldsymbol{\Sigma} = \mathbf{R}\mathbf{S}\mathbf{S}^\top\mathbf{R}^\top$. Colour may be view-dependent via spherical harmonics; the implementation described later also accepts per-Gaussian RGB in linear space.

The spatial falloff is
\[
G(\mathbf{x}) = \exp\bigl(-\tfrac{1}{2}(\mathbf{x}-\boldsymbol{\mu})^\top \boldsymbol{\Sigma}^{-1} (\mathbf{x}-\boldsymbol{\mu})\bigr).
\]

\begin{figure}[H]
\centering
\thesisfig{\textwidth}{figures/fig3_2_gaussian_splatting_projection.pdf}
\caption{Conceptual illustration of 3D Gaussian projection and splatting: each 3D Gaussian is projected onto the image plane; footprints are depth-sorted and alpha-composited \cite{kerbl20233dgs}.}
\label{fig:gaussian_project}
\end{figure}

\section{Projection, Rasterisation, and Real-Time Rendering}
\label{sec:3dgs_splatting}

Rendering splats each Gaussian onto the image \cite{kerbl20233dgs}:

\begin{enumerate}
    \item \textbf{Project} $\boldsymbol{\mu}$ with the camera matrix.
    \item \textbf{Project covariance} to $2\times2$ image-space covariance using the Jacobian $\mathbf{J}$ of the perspective mapping at $\boldsymbol{\mu}$.
    \item \textbf{Sort} primitives (globally or per tile) for approximate depth order.
    \item \textbf{Rasterise} with alpha blending, accumulating colour and transmittance along rays.
\end{enumerate}

At each pixel, ordered compositing follows
\[
\mathbf{c}_{\mathrm{out}} = \alpha \cdot \mathbf{c}_{\mathrm{G}} + (1-\alpha) \mathbf{c}_{\mathrm{in}}.
\]

\section{Differentiable Rendering and Optimisation}
\label{sec:3dgs_differentiable}

To fit Gaussians to photographs, the renderer must be \textbf{differentiable}. Gradients of a photometric loss propagate backwards through alpha compositing, rasterisation, and projection, updating $\boldsymbol{\mu}$, scales, and rotations. Optimisation typically alternates gradient steps with \textbf{adaptive density control} (splitting large Gaussians, culling transparent ones) \cite{kerbl20233dgs}. Initialisation often starts from a sparse point cloud (e.g.\ from COLMAP SfM), linking 3DGS training back to the photogrammetry chain in Chapter~\ref{ch:theory1}. Full training remains \textbf{off-device} in this application; the iOS code imports and renders only the optimised \texttt{.ply} result.

\section{Meshes versus Gaussians: Editability and Role in the Application}
\label{sec:3dgs_vs_mesh}

\textbf{Meshes (USDZ from Object Capture)} expose \textbf{explicit connectivity}: faces can be simplified or exported; axis-aligned bounding boxes align with RoomPlan furniture detections; RealityKit loads USDZ natively. \textbf{3D Gaussians} \cite{kerbl20233dgs} do not define a thin surface: coherent editing requires transforming thousands of primitives; boolean operations and CAD-style placement are ill-defined without mesh conversion.

The product decision in this thesis is therefore:

\begin{itemize}
    \item \textbf{Authoring and replacement}: USDZ meshes only (Asset~A).
    \item \textbf{Visualisation and comparison}: optional \texttt{.ply} splats via Metal (Assets~B and~C).
    \item \textbf{Desktop refinement}: Blender export carries glTF room structure and USDZ furniture copies; splat \texttt{.ply} files are not bundled in the ZIP and remain a separate in-app benchmark unless imported manually in Blender.
    \item \textbf{Training}: off-device for Asset~B; feedforward inference on a LAN Mac for Asset~C.
\end{itemize}

\begin{table}[htbp]
\centering
\caption{Mesh-based Object Capture output versus imported 3D Gaussian Splatting in this application.}
\label{tab:mesh_vs_3dgs}
\begin{tabular}{|p{0.22\textwidth}|p{0.34\textwidth}|p{0.34\textwidth}|}
\hline
\textbf{Criterion} & \textbf{Mesh (USDZ)} & \textbf{3DGS (\texttt{.ply})} \\
\hline
Geometry & Explicit triangles & Cloud of Gaussians \\
\hline
Rendering & RealityKit / Metal mesh & Splatting \cite{kerbl20233dgs} \\
\hline
RoomPlan alignment & Scale-to-bounds, grounding & Manual alignment if needed \\
\hline
Editing in app & Yes (replace, transform) & View / compare only \\
\hline
Blender export & Included (copies) & Not included in export ZIP \\
\hline
\end{tabular}
\end{table}

\section{Three Object-Level Reconstruction Paths}
\label{sec:three_assets}

Experimental comparisons in this thesis evaluate physical objects or room regions through three pipelines with different input labour, compute location, and editability. Only Asset~A integrates with interactive furniture replacement; Assets~B and~C share the Metal splat viewer.

\textbf{Asset~A (Object Capture).} Guided multi-view orbit on device; photogrammetry session produces metric textured USDZ for the furniture library and Blender export copies.

\textbf{Asset~B (offline 3DGS).} Fifty to two hundred or more photographs with poses from COLMAP-class SfM on a workstation; GPU training yields a standard \texttt{.ply} imported into the app. Highest plenoptic quality when capture discipline is good; relative scale unless metric constraints are added.

\textbf{Asset~C (SHARP).} Single JPEG posted to a user-controlled Mac on the local network; feedforward model returns a pruned \texttt{.ply} ($\leq 300\,000$ Gaussians) for mobile viewing. Minimises capture friction at the cost of single-view ambiguity and network dependency.

\begin{table}[htbp]
\centering
\caption{Three reconstruction paths compared (object-level focus).}
\label{tab:three_assets}
\small
\begin{tabular}{|p{2.1cm}|p{3.1cm}|p{3.1cm}|p{3.1cm}|}
\hline
\textbf{Criterion} & \textbf{A: Object Capture} & \textbf{B: Offline 3DGS} & \textbf{C: SHARP} \\
\hline
Input & Multi-view orbit & Multi-view + SfM poses & Single image \\
\hline
Where it runs & iPhone/iPad & External GPU & Mac on LAN \\
\hline
Time-to-result & Minutes & Minutes--hours & $\approx 10\,$s (warm server) \\
\hline
Output & USDZ mesh & \texttt{.ply} splat & \texttt{.ply} splat \\
\hline
Edit in app & Yes & View only & View only \\
\hline
Renderer & RealityKit & Metal splats & Metal splats \\
\hline
\end{tabular}
\end{table}

\section{Limitations of Classical 3DGS for Mobile Capture}
\label{sec:3dgs_limitations}

Despite its visual quality and real-time rendering speed, the classical 3DGS pipeline~\cite{kerbl20233dgs} carries practical constraints for mobile, user-initiated capture: multi-view input, offline SfM, per-scene optimisation, and typically relative scale from COLMAP initialisation. These constraints motivate feedforward SHARP (Asset~C) for rapid previews while keeping full training off-device for Asset~B.

\section{SHARP: Feedforward Single-Image 3D Gaussian Splatting}
\label{sec:sharp}

\textbf{SHARP} (\emph{Sharp Monocular View Synthesis})~\cite{sharp-monocular-view} predicts a complete Gaussian field from one RGB image in a single forward pass, without per-scene optimisation. It builds on a monocular metric depth prior, initialises on the order of $10^6$ Gaussians, and decodes attributes compatible with Kerbl et al.'s rasterisation model~\cite{kerbl20233dgs}. Table~\ref{tab:3dgs_vs_sharp} contrasts offline optimised 3DGS (Asset~B) with SHARP (Asset~C).

\begin{table}[h]
\centering
\caption{Classical 3DGS~\cite{kerbl20233dgs} versus SHARP~\cite{sharp-monocular-view}.}
\label{tab:3dgs_vs_sharp}
\setlength{\tabcolsep}{4pt}
\begin{tabular}{p{2.8cm} p{4.3cm} p{4.3cm}}
\toprule
\textbf{Aspect} & \textbf{Classical 3DGS (Asset~B)} & \textbf{SHARP (Asset~C)} \\
\midrule
Input             & Multi-view images + poses & Single image \\
Pose estimation   & COLMAP SfM (offline)      & Not required \\
Inference mode    & Per-scene optimisation    & Single forward pass \\
Inference time    & Minutes to hours          & ${<}$1\,s (GPU) / ${\approx}$10\,s (Apple Silicon) \\
Mobile suitability & Import \texttt{.ply} only & LAN server $\rightarrow$ import \texttt{.ply} \\
\bottomrule
\end{tabular}
\end{table}

\subsection{Integration in This Application}
\label{subsec:sharp_integration}

The application integrates SHARP using a client--server architecture: a Flask server (\texttt{sharp\_server.py}) on the local network exposes \texttt{/sharp/infer}; the iOS client (\texttt{SHARPService}, \texttt{SHARPCaptureView}) POSTs a JPEG with extended timeouts and displays the returned pruned \texttt{.ply} in the Metal viewer. Implementation details appear in Chapter~\ref{ch:implementation}.

This positions the application to demonstrate a reconstruction spectrum: \textbf{RoomPlan} (parametric, real-time) $\rightarrow$ \textbf{Object Capture} (photogrammetric mesh, on-device) $\rightarrow$ \textbf{offline 3DGS} (Asset~B) $\rightarrow$ \textbf{SHARP} (Asset~C) $\rightarrow$ optional \textbf{Blender export} of the staged mesh-based room (Chapter~\ref{ch:implementation}).

\chapter{Requirements and System Specification}
\label{ch:requirements}

This chapter specifies the intended behaviour of the application from a \textbf{requirements engineering} perspective: who uses the system, which scenarios must succeed, which quality attributes matter, and on which platform the software runs. Requirements are labelled for traceability to design and tests in Chapter~\ref{ch:implementation}. The specification follows the workflow \emph{capture once} $\rightarrow$ \emph{edit persistently} $\rightarrow$ \emph{optionally benchmark radiance fields} $\rightarrow$ \emph{optionally export to Blender}.

\section{Target Users and Interaction Scenarios}

\textbf{Primary user.} The end user is a person redesigning a home or office room who owns a compatible iPhone or iPad and wants to preview furniture changes without permanent purchases.

\textbf{Secondary participants.} A developer extending the codebase benefits from a clean architecture.

\textbf{Core use cases} include: (UC-1) capture a room with RoomPlan and save it; (UC-2) open the saved room in the 3D editor; (UC-3) select a detected furniture instance and assign a replacement USDZ model; (UC-4) adjust position, orientation, and scale with gestures; (UC-5) add decor (wall art, carpet, lighting) and leave the app with edits intact on return; (UC-6) view room measurements, material estimates, and rule-based Smart Tips; (UC-7a) capture a standalone object with Object Capture and add it to the furniture library (Asset~A); (UC-7b) import an offline-trained 3DGS \texttt{.ply} (Asset~B) and view it for comparison; (UC-7c) capture a single photograph and receive a SHARP \texttt{.ply} from a local server (Asset~C); (UC-8) export a staged room as a Blender-oriented ZIP (glTF, USDZ copies, manifests, helper script); (UC-9) delete a room with confirmation and artefact cleanup.

\section{Functional Requirements}
\label{sec:func_req}

\begin{description}
    \item[R-1 Room capture.] The system shall run a RoomPlan-based capture session and produce a \texttt{CapturedRoom} containing walls, floors, ceilings where detected, doors, windows, and classified objects with transforms and dimensions.

    \item[R-2 Persistence of scan data.] The system shall serialise each room's \texttt{CapturedRoom} to application-accessible storage (JSON under \texttt{Documents/RoomScans/}) keyed by room identifier, and reload it on demand.

    \item[R-3 Editor scene.] The system shall construct a three-dimensional scene in RealityKit representing the captured geometry and furniture placeholders or replacements, without requiring a live AR session for editing.

    \item[R-4 Furniture replacement.] For each detected object selected by the user, the system shall allow binding a USDZ asset (catalog or user-captured) and shall scale and position it to approximate the detected bounding volume.

    \item[R-5 Manipulation.] The user shall translate, rotate, and scale replacement models using on-screen gestures or equivalent controls; overrides shall persist in \texttt{FurnitureReplacement} records.

    \item[R-6 Object Capture.] The system shall expose a guided object-scanning flow using Object Capture APIs and shall write resulting USDZ files to a location usable by the furniture library.

    \item[R-7 Gaussian splat viewing.] The system shall parse a standard 3DGS \texttt{.ply} with expected attribute layout and render it with GPU splatting for interactive orbit viewing (Assets~B and~C).

    \item[R-8 Room deletion.] The system shall allow the user to delete any saved room from the room list, with confirmation, removing associated persistent artefacts (room index entry, scan JSON, replacement records, cached PLY files where applicable).

    \item[R-9 SHARP single-image inference.] The system shall provide a capture interface that sends a single JPEG to a user-configured SHARP inference server over the local network, receives a 3DGS \texttt{.ply}, and stores it under Documents. The HTTP client shall use a request timeout of at least 360\,s.

    \item[R-10 Server reachability and configuration.] The system shall perform a health-check ping to the configured SHARP server URL before each inference session and display reachability status. The user shall be able to configure the server base URL within the app.

    \item[R-11 Room analytics.] From a stored \texttt{CapturedRoom}, the system shall compute metric dimensions, renovation material estimates (paint, flooring, skirting heuristics), and rule-based Smart Tips including lighting density guidance derived from placed fixture count versus floor area.

    \item[R-12 Blender export.] The system shall export a saved room that has scan data as a shareable ZIP containing: (a)~a glTF scene with structural geometry and placement markers in metres; (b)~copies of USDZ furniture and wall-art images referenced by manifests; (c)~\texttt{scene\_manifest.json} for furniture and lighting placement; (d)~a Blender Python script (\texttt{blender\_apply\_furniture.py}) and README with import instructions. Export shall be available from the room list and the room designer.
\end{description}

\section{Non-Functional Requirements}
\label{sec:nfr_req}

\begin{description}
    \item[N-1 Performance.] On supported hardware, room editing shall remain interactive; splat scenes shall target $\leq 300\,000$ Gaussians at comfortable frame rates.

    \item[N-2 Privacy.] Room scanning and object reconstruction shall not require uploading scan data to a third-party server for Apple API paths; Blender export remains on-device; SHARP traffic is limited to a user-controlled LAN server.

    \item[N-3 Usability.] The user shall receive guidance during RoomPlan and Object Capture sessions. Terminology in the UI shall remain consistent with RoomPlan semantics where applicable.

    \item[N-4 Maintainability.] Application logic shall be partitioned into models, services, and SwiftUI views, following observable patterns suitable for SwiftUI.

    \item[N-5 Portability within iOS.] The project shall target a documented minimum iOS version matching RoomPlan and Object Capture availability; features degrade gracefully when optional hardware (LiDAR) is absent where Apple allows.

    \item[N-6 Network resilience.] The SHARP inference client shall use a dedicated \texttt{URLSession} with 360\,s request timeout and 600\,s resource timeout. Failures shall display descriptive errors without requiring an app restart.

    \item[N-7 Responsiveness.] Operations taking more than 0.5\,s shall display a progress indicator. The SHARP capture UI shall warn that the first inference may take approximately two minutes due to Metal shader compilation on the server.

    \item[N-8 Error handling.] The system shall handle without crashing: corrupt \texttt{.ply} import; SHARP non-200 responses; network timeout; Object Capture cancelled with insufficient coverage; delete of a room currently open in the editor; Blender export when scan data is missing.
\end{description}

\section{System Specification}
\label{sec:system_spec}

\textbf{Programming language and UI.} Swift~5, SwiftUI, observation for reactive updates.

\textbf{Graphics and AR.} RealityKit for the room editor; ARKit and RoomPlan for capture; Metal for the Gaussian splat renderer (\texttt{MTKView} bridge from SwiftUI).

\textbf{Data formats.} \texttt{CapturedRoom} via JSON (\texttt{Codable}); USDZ for meshes; PLY for splats; ZIP for Blender export (glTF, manifests, script).

\textbf{Storage layout (conceptual).} Room scan files under \texttt{Documents/RoomScans/}; USDZ, wall art, and SHARP outputs in typed subdirectories; lightweight room metadata in UserDefaults.

\begin{table}[htbp]
\centering
\caption{Requirement identifiers and short descriptions (traceability).}
\label{tab:req_trace}
\setlength{\tabcolsep}{6pt}
\renewcommand{\arraystretch}{1.25}
\begin{tabular}{lp{0.68\textwidth}}
\toprule
\textbf{ID} & \textbf{Summary} \\
\midrule
\rowcolor{gray!8}
R-1 & RoomPlan capture session $\rightarrow$ \texttt{CapturedRoom} \\
R-2 & Serialise / reload room scan JSON per room identifier \\
\rowcolor{gray!8}
R-3 & RealityKit non-AR editor scene with orbit camera \\
R-4 & USDZ asset binding, scaling, and placement per detected object \\
\rowcolor{gray!8}
R-5 & Gesture-based translate / rotate / scale with persistence \\
R-6 & Object Capture $\rightarrow$ USDZ (Asset~A) \\
\rowcolor{gray!8}
R-7 & Binary PLY import and Metal GPU splatting (Assets B/C) \\
R-8 & Room deletion with confirmation and artefact cleanup \\
\rowcolor{gray!8}
R-9 & SHARP single-image inference via local server \\
R-10 & Server health-check; URL configuration \\
\rowcolor{gray!8}
R-11 & Measurements, material estimates, Smart Tips \\
R-12 & Blender export ZIP (glTF + USDZ + script) \\
\midrule
\rowcolor{gray!8}
N-1--N-8 & Performance, privacy, usability, maintainability, portability, network, responsiveness, errors \\
\bottomrule
\end{tabular}
\end{table}

\begin{figure}[H]
\centering
\thesisfig{0.85\textwidth}{figures/ch4_context_diagram.pdf}
\caption{System context: user, application, and main persistent artefacts (room JSON, USDZ, PLY, Blender export ZIP). SHARP runs on a user-controlled Mac on the local network.}
\label{fig:context_req}
\end{figure}

\section{Constraints and Dependencies}

The application depends on Apple's RoomPlan and Object Capture frameworks \cite{apple2024roomplan2}. 3DGS quality for Asset~B depends on external training discipline. SHARP requires a user-maintained server. Blender export requires persisted \texttt{CapturedRoom} data. No requirement mandates a specific commercial furniture marketplace.

\section{Risk Analysis}

Table~\ref{tab:risks} identifies principal technical risks and mitigations.

\begin{table}[h]
\centering
\caption{Risk analysis: risks, likelihood, impact, and mitigation.}
\label{tab:risks}
\setlength{\tabcolsep}{6pt}
\renewcommand{\arraystretch}{1.3}
\begin{tabular}{p{3.0cm} c c p{5.2cm}}
\toprule
\textbf{Risk} & \textbf{Likelihood} & \textbf{Impact} & \textbf{Mitigation} \\
\midrule
\rowcolor{gray!8}
SHARP server unavailable & Medium & Medium & Health-check; graceful error; other features independent \\
SHARP first-run timeout & High & Medium & 360\,s timeout; UI warning \\
\rowcolor{gray!8}
Large \texttt{.ply} & High & High & Server opacity prune $\leq 300\,k$ Gaussians \\
RoomPlan incomplete on non-LiDAR & Medium & High & Feature gating where applicable \\
\rowcolor{gray!8}
Blender export without scan & Medium & Medium & Disable export; alert user \\
Accidental room deletion & Low & High & Confirmation alert; cancel default \\
\bottomrule
\end{tabular}
\end{table}

\chapter{Design, Implementation, and Testing}
\label{ch:implementation}

This chapter describes how the requirements of Chapter~\ref{ch:requirements} are realised in software. It presents the \textbf{architectural decomposition}, the \textbf{data path} from RoomPlan scan to persisted room and RealityKit entities, the roles of major Swift types and services, the \textbf{user interface} structure (including room analytics and Blender export), and a \textbf{testing} approach suitable for a bachelor project.

\section{Architectural Overview}

The application entry point is a SwiftUI \texttt{App} type that configures core services and hosts a root view responsible for the local interior-design workflow, including room scanning, editing, and object capture.

Architecturally, the code follows patterns compatible with \textbf{SwiftUI and MVVM}: views observe reference-type \texttt{ObservableObject} managers (\texttt{RoomManager}, \texttt{FurnitureCatalogManager}, \texttt{ScannedFurnitureManager}, \texttt{ProjectManager}) injected through the environment. Long-running capture sessions (RoomPlan, Object Capture) live in dedicated views that bridge UIKit where needed. Rendering uses \textbf{RealityKit} (\texttt{ARView} in non-AR camera mode) for the room editor and \textbf{Metal} via \texttt{MTKView} for Gaussian splats.

The architecture comprises four layers (Figure~\ref{fig:architecture}):

\begin{enumerate}
    \item \textbf{SwiftUI Views} --- \texttt{MainTabBar}, \texttt{RoomListView}, \texttt{RoomDesignerView}, \texttt{RoomScanView}, \texttt{ObjectScanView}, \texttt{GaussianSplatView}, \texttt{SHARPCaptureView}, \texttt{RoomStatisticsView}, and related screens.
    \item \textbf{Observable Managers} --- domain state and persistence triggers.
    \item \textbf{Services} --- \texttt{SHARPService}, \texttt{GaussianSplatRenderer}, \texttt{MeasurementService}, \texttt{SmartSuggestionsService}, \texttt{RoomColorCaptureService}, \texttt{RoomBlenderExportService}, and others.
    \item \textbf{Apple Frameworks and Storage} --- RoomPlan, Object Capture, RealityKit, Metal; artefacts in the Documents directory.
\end{enumerate}

\begin{figure}[htbp]
\centering
\thesisfig{0.95\textwidth}{figures/ch5_arch_layers.png}
\caption{Layered application architecture: SwiftUI views observe injectable managers, which delegate to services and Apple frameworks. Storage artefacts include room JSON, USDZ, PLY, and Blender export ZIPs.}
\label{fig:architecture}
\end{figure}

\section{Data Model and Persistence}
\label{sec:ch5_data_model}

The domain type \texttt{Room} (\texttt{Codable}) aggregates identity, timestamps, \texttt{FurnitureReplacement} records, optional \texttt{RoomMaterials}, and optional \texttt{projectId}. It does \textbf{not} embed the full \texttt{CapturedRoom} in the room-list JSON; scan data is stored separately under \texttt{Documents/RoomScans/\{uuid\}.json} with a static in-memory cache and atomic writes. \texttt{RoomManager} owns the saved-rooms index; \texttt{cleanOrphanedScanFiles} removes JSON after room deletion (R-8).

\texttt{FurnitureReplacement} supports USDZ swaps and virtual decor (wall art, carpet, lighting), discriminated by \texttt{originalCategory}, persisted on the same path as furniture replacements. Figures~\ref{fig:domain_model} and~\ref{fig:storage_layout} summarise the model and on-disk layout.

\begin{figure}[htbp]
\centering
\thesisfig{0.97\textwidth}{figures/ch5_domain_model.png}
\caption{Domain model: \texttt{Room} aggregates replacements; \texttt{CapturedRoom} is a sidecar JSON file.}
\label{fig:domain_model}
\end{figure}

\begin{figure}[htbp]
\centering
\thesisfig{0.88\textwidth}{figures/ch5_storage_layout.png}
\caption{On-device storage layout under the app sandbox.}
\label{fig:storage_layout}
\end{figure}

\section{Data Flow: Scan to Editor}

Figure~\ref{fig:dataflow} illustrates capture $\rightarrow$ persistence $\rightarrow$ editor $\rightarrow$ relaunch. \texttt{RoomScanView} finalises \texttt{RoomCaptureSession.stop(pauseARSession: true)}; \texttt{RoomManager} writes metadata and scan JSON; \texttt{RoomDesignerView} loads the scan, builds RealityKit entities, and writes back gesture and replacement state on save.

\begin{figure}[htbp]
\centering
\thesisfig{0.97\textwidth}{figures/ch5_dataflow.png}
\caption{Data flow from RoomPlan capture through persistence to the RealityKit editor.}
\label{fig:dataflow}
\end{figure}

\section{Key Modules and Views}

\texttt{MainTabBar} separates \textbf{Furniture} (library, Object Capture, splat tools), \textbf{Rooms} (list, scan, editor, Blender export from context menu), and \textbf{Projects} (multi-room grouping). \texttt{RoomDesignerView} hosts the non-AR \texttt{ARView}, furniture picking, materials, statistics sheet, and toolbar export for Blender (R-12). \texttt{ObjectScanView} registers completed USDZ in \texttt{ScannedFurnitureManager}. The splat pipeline (\texttt{GaussianSplatData}, \texttt{GaussianSplatRenderer}, \texttt{GaussianSplatView}) implements R-7.

\subsection{SHARP single-image inference (R-9, R-10)}
\label{sec:sharp}

\texttt{SHARPService} POSTs JPEG data to \texttt{/sharp/infer} using a dedicated \texttt{URLSession} (360\,s request, 600\,s resource timeout). The server (\texttt{sharp\_server.py}) prunes Gaussians by opacity to $\leq 300\,000$ before returning the \texttt{.ply}. A \texttt{GET /health} ping with 4\,s timeout precedes capture (R-10). Figure~\ref{fig:sharp_pipeline} summarises the path.

\begin{figure}[htbp]
\centering
\thesisfig{0.97\textwidth}{figures/ch5_sharp_pipeline.png}
\caption{SHARP inference pipeline: device $\rightarrow$ LAN server $\rightarrow$ \texttt{.ply} $\rightarrow$ Metal viewer.}
\label{fig:sharp_pipeline}
\end{figure}

\section{USDZ Alignment Heuristics}
\label{sec:alignment}

When a replacement is first bound, uniform scale $s = \min(d_x^{\mathrm{room}}/d_x^{\mathrm{usdz}}, d_y^{\mathrm{room}}/d_y^{\mathrm{usdz}}, d_z^{\mathrm{room}}/d_z^{\mathrm{usdz}})$ fits the model within the detected bounds, followed by floor grounding from the \texttt{CapturedRoom.Object} transform. \texttt{customScale} stores user pinch adjustments.

\section{Metal Gaussian Splatting}
\label{sec:splat_rendering}

The splat renderer follows Kerbl et al.\ \cite{kerbl20233dgs}: 3D covariance from quaternion and scale, projection to $2\times2$ screen-space covariance via Jacobian $\mathbf{J}$, per-frame depth sort, and alpha compositing in fragment shaders. A low-pass filter on the 2D covariance prevents sub-pixel flicker. Budgets up to 300\,k Gaussians target interactive frame rates (N-1). Figure~\ref{fig:splat_pipeline} shows the per-frame pipeline.

\begin{figure}[htbp]
\centering
\thesisfig{0.55\textwidth}{figures/ch5_splat_pipeline.png}
\caption{Per-frame Metal Gaussian splatting pipeline.}
\label{fig:splat_pipeline}
\end{figure}

\section{Room Analytics (R-11)}
\label{sec:analytics}

\texttt{RoomStatisticsView} presents four tabs driven by services on the stored \texttt{CapturedRoom}:

\begin{itemize}
    \item \textbf{Measurements.} \texttt{MeasurementService} sums floor and wall areas, computes an axis-aligned bounding extent, perimeter, and approximate volume.
    \item \textbf{Objects.} Counts detections versus replacements and decor.
    \item \textbf{Materials.} Heuristic paint litres (assuming $10\,\mathrm{m}^2/\mathrm{L}$ and two coats after subtracting nominal opening areas), flooring area with waste factor, and skirting length.
    \item \textbf{Smart Tips.} \texttt{SmartSuggestionsService.generateSuggestions} applies rule-based checks: lighting fixture count compared to $N_{\mathrm{rec}} = \max(1, \lfloor A_{\mathrm{floor}} / 4\,\mathrm{m}^2 \rfloor)$; bare-wall and carpet coverage heuristics; furniture density thresholds. No luminance simulation is performed---only placed fixture count versus floor area.
\end{itemize}

\section{Blender Export (R-12)}
\label{sec:blender_export}

\texttt{RoomBlenderExportService.makeExportZip} assembles a temporary ZIP archive for import into Blender or other DCC tools. Export requires persisted scan data (\texttt{CapturedRoom}); otherwise the UI shows an error.

The archive typically contains:

\begin{itemize}
    \item \texttt{room\_scene.gltf} and \texttt{room\_scene.bin} --- structural boxes and green placement markers in metres (\texttt{GltfMiniBuilder}).
    \item \texttt{furniture/*.usdz} --- copies of staged assets.
    \item \texttt{wall\_art/*.jpg} --- exported wall-art textures where applicable.
    \item \texttt{furniture\_manifest.json}, \texttt{scene\_manifest.json} --- transforms, categories, lighting entries.
    \item \texttt{blender\_apply\_furniture.py} --- post-import script: import USDZ onto markers, create lights from manifest, remove guide geometry.
    \item \texttt{README.txt} --- import and script execution steps.
\end{itemize}

Export is exposed from \texttt{RoomDesignerView} (toolbar \emph{Export for Blender}) and \texttt{RoomListView} (context menu). \texttt{materialsForListExport} resolves materials when export is initiated from the list without an open designer session. The share sheet hands the ZIP to Files or macOS. Splats (\texttt{.ply}) are \textbf{not} included; they remain in-app benchmarks unless the user imports them separately in Blender.

\section{On-Device versus Off-Device Processing}

RoomPlan, Object Capture, editor persistence, analytics, and Blender ZIP creation run \textbf{on-device}. Asset~B training runs \textbf{off-device}; SHARP inference runs on a \textbf{user-controlled Mac} on the LAN (N-2).

\section{User Interface Walkthrough}

\textbf{Scanning.} The user starts RoomPlan from the rooms list; structural elements and furniture appear incrementally. On \emph{Done}, \texttt{RoomCaptureSession.stop} yields \texttt{CapturedRoom}; \texttt{RoomColorCaptureService} may sample dominant wall and floor colours during the walk-through for \texttt{RoomMaterials}.

\textbf{Editor.} \texttt{RoomDesignerView} uses non-AR orbit navigation. The user selects detections, binds USDZ, adjusts via gestures and axis gizmo, and may open \texttt{RoomStatisticsView}. Decor pickers add wall art, carpet, and lighting.

\textbf{Benchmarks (optional).} Object Capture (Asset~A), imported \texttt{.ply} (Asset~B), or SHARP capture (Asset~C) open in the splat viewer or library flows for comparison with the mesh editor.

\textbf{Export (optional).} Blender export packages the staged semantic scene for desktop refinement (R-12).

\begin{figure}[htbp]
\centering
\thesisfig{0.72\textwidth}{figures/ch5_ui_03_room_scan.png}\\[0.6em]
\thesisfig{0.72\textwidth}{figures/ch5_ui_04_editor.png}
\caption{Example screenshots: RoomPlan capture (top) and non-AR room editor (bottom).}
\label{fig:ui_scan_editor}
\end{figure}

\subsection{Demonstration checklist}

\begin{enumerate}
    \item Launch; verify saved rooms and new scan entry.
    \item Complete RoomPlan; confirm scan JSON under \texttt{RoomScans/} (R-1, R-2).
    \item Editor: replace furniture; gesture persistence after relaunch (R-3--R-5).
    \item Decor and statistics tabs; add lighting; observe Smart Tips update (R-11).
    \item Object Capture (optional) (R-6).
    \item SHARP or PLY import (optional) (R-7, R-9, R-10).
    \item Blender export; import \texttt{room\_scene.gltf} in Blender; run helper script (R-12).
    \item Delete room with confirmation (R-8).
\end{enumerate}

\section{Testing}

Manual test cases map requirements to outcomes (Table~\ref{tab:test_map}). Performance is assessed informally; splat scenes target 60\,fps at $\leq 300\,k$ Gaussians (N-1). Edge cases include corrupt JSON, oversized PLY, SHARP timeout, export without scan data, and delete-while-open (N-8).

\begin{table}[htbp]
\centering
\caption{Requirement-to-test mapping.}
\label{tab:test_map}
\begin{tabular}{|c|p{0.68\textwidth}|}
\hline
\textbf{Req.} & \textbf{Test procedure} \\
\hline
R-1, R-2 & Scan room; verify JSON in \texttt{RoomScans/}; relaunch. \\
\hline
R-3--R-5 & Editor; bind USDZ; transform; kill app; reopen. \\
\hline
R-6 & Complete object scan; USDZ in picker. \\
\hline
R-7 & Import known-good \texttt{.ply}; orbit smoothly. \\
\hline
R-8 & Delete room; confirm; scan file removed. \\
\hline
R-9, R-10 & SHARP URL; health indicator; capture; view splat. \\
\hline
R-11 & Statistics tabs; lighting tips respond to fixture count. \\
\hline
R-12 & Export ZIP; Blender import glTF; run \texttt{blender\_apply\_furniture.py}. \\
\hline
\end{tabular}
\end{table}

\section{Summary}

The implementation binds Apple's capture APIs to a persistent, editable RealityKit scene graph through explicit \texttt{CapturedRoom} storage and replacement metadata. Optional Metal splatting and SHARP provide radiance-field benchmarks without replacing the mesh-based staging model. Room analytics and Blender export extend the workflow toward planning aids and professional DCC refinement. Requirements R-1--R-12 and non-functional goals N-1--N-8 are addressed in the design decisions documented above.

\chapter{Conclusions}
\label{ch:conclusions}

This thesis investigated how a commodity iOS device can support \textbf{interior design preview} through \textbf{parametric room capture}, \textbf{mesh-based editing}, optional \textbf{plenoptic visualisation}, \textbf{rule-based room analytics}, and \textbf{Blender export} for desktop refinement. The theoretical chapters argued that Apple's Object Capture instantiates a classical \textbf{photogrammetric} chain, whereas RoomPlan builds on \textbf{depth fusion} and learning to emit a structured \texttt{CapturedRoom}. The second theory chapter positioned \textbf{3D Gaussian Splatting} \cite{kerbl20233dgs} and feedforward \textbf{SHARP} as real-time alternatives to meshes for \emph{viewing}, not for CAD-style furniture replacement, and formalised a three-path comparison (Assets~A, B, and~C).

\section{Summary of Results}

The software artefact demonstrates a coherent pipeline: \textbf{RoomPlan} $\rightarrow$ serialised scan $\rightarrow$ \textbf{RealityKit} editor $\rightarrow$ USDZ replacements with gestures and decor $\rightarrow$ persistence across sessions. \textbf{Object Capture} extends asset authoring (Asset~A). A \textbf{Metal} splat renderer implements R-7 for imported \texttt{.ply} data (Asset~B). The \textbf{SHARP} integration (R-9, R-10) provides Asset~C via a local-network inference server without third-party cloud dependency for Apple capture paths. \textbf{Room analytics} (R-11) translate scan geometry into measurements, material estimates, and heuristic Smart Tips. \textbf{Blender export} (R-12) packages glTF structure, USDZ furniture, manifests, and a helper script so users can continue staging in desktop tools.

Requirements R-1--R-12 and associated non-functional goals N-1--N-8 are addressed at the design level in Chapter~\ref{ch:requirements} and at the code level in Chapter~\ref{ch:implementation}.

\section{Comparison of Reconstruction Modes}

Qualitative evaluation during development indicated that Asset~A (Object Capture meshes) delivers manipulable geometry aligned to RoomPlan bounding volumes and supports both in-app replacement and Blender export. Asset~B (offline-trained splats) achieves strong view-dependent realism when many calibrated views are available but does not support direct editing in the app. Asset~C (SHARP) minimises capture friction (single photograph, LAN inference after server warm-up) at the cost of single-view ambiguity and network dependency. None of the splat paths automatically register to \texttt{CapturedRoom} semantics; unified AR compositing would require future registration or fusion work.

\section{Limitations}

Accuracy is bounded by RoomPlan detection quality, photogrammetry assumptions in Object Capture, and the lack of closed-loop feedback between splats and parametric room semantics. Smart Tips are illustrative heuristics, not certified lighting design. Blender export reproduces the staged semantic scene, not photorealistic radiance fields; splats are excluded from the ZIP. 3DGS training remains external; real-time splat \emph{training} on device is research-grade \cite{Guo26}. Usability has not been quantified with SUS/HARUS in this thesis text; Kandil et al.\ \cite{kandil2021ux} provide a methodology for a future study.

\section{Future Work}

Promising extensions include: formal user evaluation; tighter integration between exported Object Capture image sets and external COLMAP/3DGS pipelines; coarse automatic registration of splats to room coordinates; optional inclusion of splat assets in DCC workflows; level-of-detail for very large \texttt{.ply} files; and semantic layout suggestions informed by work such as K\'an et al.\ \cite{kan2021ar}. Mithun et al.\ \cite{Geo2023} demonstrate that synthetic datasets with rich semantic labels can improve indoor scene understanding for AR applications, complementary to the real-scan approach taken here.

The project shows that \textbf{merging API-level room semantics with an editable mesh workflow and optional Blender hand-off} yields a practical middle ground between live AR furniture preview and professional BIM tools, suitable for iterative interior exploration on mobile hardware.


\bibliography{references}

\end{document}
